\section{Formal Reasoning as a Special Case}
\label{sec:formal-reasoning}

Classical first-order logic and probabilistic graphical models are usually
presented as separate formalisms. The QBBN~\cite{coppola2024qbbn} unifies
these: it is a probabilistic graphical model over propositions with AND/OR
boolean structure, proven consistent and complete with respect to the
first-order calculus. Logical deduction is the special case where the OR
gate weights are deterministic; statistical inference is the general case
where the weights are soft and learned from data.

\subsection*{The Controlled/Learned Spectrum}

This suggests a spectrum between controlled and learned reasoning:

\begin{itemize}
\item \textbf{Fully controlled}: OR gate weights are fixed to be
deterministic. The system implements exact logical deduction. This is the
QBBN in symbolic mode.

\item \textbf{Fully learned}: OR gate weights are learned from data. The
system implements soft probabilistic inference. This is the standard LLM.

\item \textbf{Hybrid}: some weights are fixed (known logical rules) and
some are learned (uncertain statistical relationships). This is the target
architecture for systems combining flexibility with formal reliability.
\end{itemize}

The transformer's BP structure is compatible with all three modes. Current
training produces the fully learned mode.

\subsection*{The Completeness Connection --- and Its Limits}

The QBBN paper proves consistency and completeness with respect to the
first-order calculus. This is a result about \emph{representational
expressivity}, not about practical inference capability. It says the QBBN
can in principle express any first-order fact --- not that a transformer
will reliably compute it, nor that inference will be tractable, nor that
the right weights will be learned from data.

Concretely: a transformer with grounded BP weights is not thereby guaranteed
to do mathematics reliably. Mathematical reasoning requires not just the
right representational structure but also the right weights, sufficient
layers for the required inference depth, and a training procedure that
produces calibrated beliefs. The completeness result says the architecture
is not the bottleneck. It does not say the other bottlenecks are easy.

\subsection*{Fast vs.\ Slow Reasoning}

The QBBN analysis connects to Kahneman's fast/slow
distinction~\cite{kahneman2011thinking}. Fast reasoning is forward
inference: causes propagate to effects in one pass of BP. Slow reasoning
requires reasoning by cases --- exploring multiple hypotheses, maintaining
multiple belief states, backtracking. The number of transformer layers
bounds the depth of fast reasoning; slow reasoning requires iteration or
explicit search outside the model.

\subsection*{Epistemic Status}

\begin{center}
\begin{tabular}{ll}
\toprule
Claim & Status \\
\midrule
QBBN is consistent and complete w.r.t.\ first-order logic & Formal theorem \\
Transformers are compatible with formal reasoning & Direct consequence \\
Trained transformers reliably do formal reasoning & Not established \\
Completeness implies practical mathematical competence & Explicitly not claimed \\
Hybrid architecture is achievable & Architectural proposal \\
\bottomrule
\end{tabular}
\end{center}